RLHF-trained language models exhibit systematic calibration inversion---confidently predicting wrong answers---on specific benchmark task clusters. We investigate this phenomenon through calibration clustering and mechanism verification. Clustering 2,212 tasks by calibration patterns yields silhouette score 0.6016, revealing non-random failure structure. We trace this to reward signal conflation: annotators rate correctness and user-modeling tasks identically (rate\_diff = 0.001), reward models show similar confidence across task types (overlap = 0.647), and model hidden states fail to separate them (separation = 0.024). Four of five mechanism hypotheses pass validation. However, keyword-based bidirectional feature detection achieves only 2.3\% prevalence ($r = -0.027$), leaving the feature-cluster link unverified. Our findings establish calibration clustering as a diagnostic tool for RLHF behavioral analysis and identify annotator conflation as the mechanism root cause, while highlighting semantic detection as needed future work.
